\ItemCategory{}
\ItemSubCategory{}
\ItemFolder{}

\chapter*{Emberbound Fragment}\stepcounter{section}\phantomsection\addcontentsline{toc}{section}{Emberbound Fragment}\label{mi:EmberboundFragment}
\MagicItemTypeAndImage[%
	top-margin		= -1.25cm,
	bottom-margin	= -1.75cm,
	image-width		= \columnwidth,
%	image-position	= 0cm,
]{Wondrous Item, Rare (requires attunement)}{images/Magic_Items/Emberbound_Fragment.png}

%\itemDescriptionAndImage[\columnwidth]{Wondrous Item, Rare (requires attunement)}{images/Magic_Items/Emberbound_Fragment.png}{5cm}\\

\section*{Appearance}
{\entryfont This jagged fragment of blackened plate armour is all that remains of a once finely wrought suit. Its scorched metal is traced with delicate ornamentation, now broken by deep crimson cracks that seem to smoulder without giving off any warmth. Despite its ruined state, the fragment feels unnaturally heavy and clings strangely to those who handle it.}

\section*{History}
{\entryfont The Emberbound Fragment is all that remains of a once finely crafted suit of plate armour. Its creator, a sorcerer experimenting with elemental binding, sought to imprison a Fire Elemental within the armour and harness its power. Before the ritual could be completed, however, the sorcerer was interrupted and the armour was shattered.
The unfinished binding survived its destruction. A portion of the elemental remained trapped within one of the fragments, unable to escape and constantly starved of the elemental energy meant to sustain it. It now draws Fire energy from its surroundings and from whoever becomes bound to the fragment, growing increasingly active as its hunger is fed.}

\section*{Magic}
{\entryfont When touching the fragment, the creature immediately becomes attuned to it. This attunement ignores the normal time required and does not require any consent.

\paragraph*{Curse} This item is cursed. Attuning to it curses you until you are targeted by a Remove Curse spell or similar magic. As long as you remain cursed, you are unwilling to part with the fragment. If it leaves your possession, it magically reappears in your pocket, pack, or among your belongings at the next dawn.

The curse deepens the longer you remain attuned to the fragment. The later stages of this curse cannot be detected, identified, or revealed by spells such as identify, divination magic, experimentation, or any other means before they manifest. When a new stage takes effect, its effects become immediately obvious to the attuned creature.}
\subsection*{Stage 1: Smothered Embers}
{\entryfont After four days attuned to the fragment, whenever you make an Attack Roll for a Spell or Weapon Attack that deals Fire damage, you take a -1 penalty to the Attack Roll. In addition, whenever you roll Fire Damage, reduce the result of each individual Damage Die by 1 (minimum of 1 per die).}
\subsection*{Stage 2: Fading Inferno}
{\entryfont After seven days attuned to the fragment, the penalty to Attack Rolls becomes -2, and each individual Fire Damage Die is reduced by 2 (minimum of 1 per die).}
\subsection*{Stage 3: Searing Reprisal}
{\entryfont After nine days attuned to the fragment, you gain \textbf{Resistance to Fire Damage}. Whenever this Resistance reduces Fire Damage you take, the fragment absorbs an amount of Fire Damage equal to the damage prevented. In addition, whenever the penalty imposed by Fading Inferno reduces Fire Damage you deal, the fragment absorbs an amount equal to the damage prevented.

On Initiative Count 20 (losing initiative ties), the stored energy erupts from it. You and each other creature in a 10-foot-radius Sphere centred on you must make a DC 16 Dexterity Saving Throw, taking 1d6 Fire Damage for every 4 Fire Damage stored in the fragment (rounded down) on a failed save, or half as much damage on a successful one. The fragment then loses all stored Fire Damage.

Fire Damage caused by this eruption cannot be reduced by the Resistance to Fire Damage granted by the fragment.}

\subsection*{Lessen the Effects}
{\entryfont Knowledgeable individuals throughout Fife may be able to suppress some of its effects or slow its progression. If the characters seek such assistance, see the \textbf{Emberbound Curse} sections in \hyperref[subsec:AuchtertoolEmberboundCurse]{\LinkFont{Auchtertool}} (p. \pageref{subsec:AuchtertoolEmberboundCurse}), in \hyperref[subsec:MethvenEmberboundCurse]{\LinkFont{Methven}} (p. \pageref{subsec:MethvenEmberboundCurse}), and in \hyperref[subsec:WoodsOfLomondEmberboundCurse]{\LinkFont{Woods of Lomond}} (p. \pageref{subsec:WoodsOfLomondEmberboundCurse}).}

\subsection*{Removing the Curse}
{\entryfont When Remove Curse or similar magic is cast on the fragment, the trapped elemental violently manifests to resist the spell. The manifestation appears in the nearest unoccupied space to the spellcaster and is hostile, focusing its attacks on the spellcaster whenever possible.

The creature that manifests depends on the current stage of the curse: a \hyperref[monster:Firegeist]{\LinkFont{Firegeist}} at Stage 1, a \hyperref[monster:FireElemental]{\LinkFont{Fire Elemental}} at Stage 2, or a \hyperref[monster:FireElementalMyrmidon]{\LinkFont{Fire Elemental Myrmidon}} at Stage 3.

If the manifestation is defeated, the elemental bound to the fragment is destroyed, immediately ending the curse and the creature's attunement to the fragment. The Broken Armour Piece then becomes mundane. If the spellcaster is reduced to 0 Hit Points before the manifestation is defeated, the manifestation is immediately drawn back into the fragment and the curse remains unaffected.

For an alternative means of resolving the curse, see [Quest Name] (p. XX).}